\documentclass[11pt]{article}
\usepackage[margin=1in]{geometry}
\usepackage{amsmath,amssymb,booktabs,array}
\usepackage[hidelinks]{hyperref}

\title{MAH Model Rebuild: Main-Text Model Candidate}
\author{}
\date{Phase 16}

\begin{document}
\maketitle

\begin{abstract}
This candidate is the restrained main-text split of the approved Phase 15
derivation draft. The MAH reform makes retained entrusted manufacturing
legally available by changing one institutional wedge. Developers choose one
common project-advancement intensity before project characteristics and
organizational routes are realized. Route choice is deterministic, and the
price of qualified manufacturing services is endogenous. The baseline
contains neither patent-generating research nor a fixed ranking between
original and incremental project responses.
\end{abstract}

\setcounter{section}{2}
\section{Model}
\label{sec:p16m-model}

\subsection{Environment and timing}
\label{subsec:p16m-environment}

There is a continuum of developers. Developer \(i\) has predetermined
research capability \(a_i>0\) and internal manufacturing capability
\(k_i>0\), distributed according to \(H(a,k)\). A viable project that reaches
commercialization-relevant planning draws value \(q>0\) and manufacturing
requirement \(m>0\) from the exogenous distribution \(F(q,m)\). For empirical
decompositions only,
\[
  F(q,m)=\sum_{g\in\{O,\mathrm{Inc}\}}\rho_gF_g(q,m),
  \qquad
  \rho_g\geq0,\quad \sum_g\rho_g=1.
\]
The class label \(g\) creates no class-specific control.

The institutional regime is \(M\in\{0,1\}\). The pre-MAH regime sets the
barrier to retained entrusted manufacturing at infinity; the MAH regime
makes it finite:
\begin{equation}
  \tau_E(0)=+\infty,
  \qquad
  \tau_E(1)=\bar\tau_E<+\infty,
  \qquad \bar\tau_E\geq0.
  \label{eq:p16m-policy}
\end{equation}
The available routes are internal manufacturing \(I\), retained entrusted
manufacturing \(E\), non-retained transfer \(T\), and abandonment or delay
\(A\). Under \(E\), the developer remains the authorization holder and a
qualified external producer manufactures. Thus \(E\) is not transfer.

Before project characteristics and routes are realized, developer \(i\)
chooses one common \(x_i\geq0\), interpreted as original-drug innovation
investment / project-advancement intensity. It is broader than pure clinical
development, but excludes basic research, compound discovery,
patent-generating effort, and patent applications. The expected measure of
viable projects reaching route planning is
\begin{equation}
  \lambda_i^{\mathrm{plan}}=a_ix_i.
  \label{eq:p16m-planning}
\end{equation}

Timing is as follows. First, \(M\) is observed and developers anticipate the
market-consistent qualified-service price. Second, they choose \(x_i\).
Third, planning-stage projects draw \((q,m)\) and choose a route. Fourth,
downstream realization occurs with exogenous probability \(s(q)\). Finally,
aggregate capacity demand and supply must validate the anticipated service
price. The last step is an equilibrium consistency condition, not a late
unanticipated shock. Upstream research and patent generation may help
determine the predetermined \(a_i\), but they are outside the baseline
decision problem.

\subsection{Commercialization technologies}
\label{subsec:p16m-technologies}

Conditional on successful commercialization, residual demand is
\(y(p;q)=Aqp^{-\varepsilon}\), where \(A>0\) and \(\varepsilon>1\).
Product-price optimization gives
\[
  p^*(c)=\frac{\varepsilon}{\varepsilon-1}c
\]
and the gross operating present value
\begin{equation}
  R(q,c)
  =
  \frac{Aq}{1-\beta\varphi}
  \frac{(\varepsilon-1)^{\varepsilon-1}}
       {\varepsilon^\varepsilon}
  c^{1-\varepsilon},
  \qquad
  \beta\in(0,1),\quad \varphi\in[0,1].
  \label{eq:p16m-return}
\end{equation}
The appendix derives the FOC, SOC, and accounting. In particular,
\(R_q>0\) and \(R_c<0\).

Internal production uses marginal cost \(c_I(m,k_i)\), with
\(c_{I,m}>0\) and \(c_{I,k}<0\), and setup cost \(F_I(m,k_i)\), with
\(F_{I,m}>0\) and \(F_{I,k}<0\) on the feasible domain. Internal production
is unavailable when \(k_i<\underline{k}(m)\), represented by
\(F_I(m,k_i)=+\infty\). Its route value is
\begin{equation}
  W_i^I(q,m)
  =
  s(q)R\!\left(q,c_I(m,k_i)\right)-F_I(m,k_i).
  \label{eq:p16m-internal}
\end{equation}

Entrusted production uses external marginal cost \(c_E(m)\), requires
\(b(m)>0\) units of qualified capacity with \(b'(m)>0\), incurs setup cost
\(F_E(m)\), and leaves the holder with burden \(\mu_E\geq0\). At the
qualified-capacity price \(p_m\),
\begin{equation}
  \begin{aligned}
  W_i^E(q,m;M,p_m)
  ={}&
  s(q)R\!\left(q,c_E(m)\right)-F_E(m)-p_mb(m)\\
  &-\mu_E-\tau_E(M).
  \end{aligned}
  \label{eq:p16m-entrusted}
\end{equation}
Each monetary cost enters once. The transfer and abandonment values are
\begin{equation}
  W^T(q,m)=T(q,m),
  \qquad
  W^A=0.
  \label{eq:p16m-outside}
\end{equation}

\subsection{Organization and project advancement}
\label{subsec:p16m-organization}

At a fixed \(p_m\), the project chooses deterministically:
\begin{equation}
  \begin{aligned}
  W_i(q,m;M,p_m)
  &=
  \max\{W_i^I,W_i^E,W^T,W^A\},\\
  r_i^*(q,m;M,p_m)
  &\in
  \arg\max_{r\in\{I,E,T,A\}} W_i^r.
  \end{aligned}
  \label{eq:p16m-route-choice}
\end{equation}
Continuous heterogeneity makes ties a zero-measure event; no logit or
inclusive value is used.

The internal-minus-entrusted gap is
\begin{equation}
  \begin{aligned}
  \Delta_{IE}(k_i)
  ={}&
  s(q)\!\left[
  R\!\left(q,c_I(m,k_i)\right)-R\!\left(q,c_E(m)\right)
  \right]-F_I(m,k_i)\\
  &+F_E(m)+p_mb(m)+\mu_E+\tau_E(M).
  \end{aligned}
  \label{eq:p16m-gap}
\end{equation}
Holding \((q,m,M,p_m)\) fixed on the internally feasible domain,
\[
  \Delta_{IE,k}
  =
  s(q)R_c\!\left(q,c_I\right)c_{I,k}-F_{I,k}>0.
\]
When the endpoint crossing conditions hold, a unique conditional cutoff
\(k^*(q,m;p_m,M)\) satisfies \(\Delta_{IE}(k^*)=0\). Conditional on \(I\)
and \(E\) dominating \(T\) and \(A\), developers below the cutoff select
\(E\), while those above it select \(I\). Transfer and abandonment remain
valid outside options.

Expected optimized value per planning-stage project is
\begin{equation}
  \Omega_i(M,p_m)
  =
  \int W_i(q,m;M,p_m)\,dF(q,m).
  \label{eq:p16m-omega}
\end{equation}
Developer \(i\) solves
\begin{equation}
  \max_{x_i\geq0}
  \left\{
  \beta a_ix_i\Omega_i(M,p_m)
  -
  \frac{\kappa}{1+\nu}x_i^{1+\nu}
  \right\},
  \qquad \kappa>0,\quad \nu>0.
  \label{eq:p16m-advancement-problem}
\end{equation}
Strict concavity gives the unique solution
\begin{equation}
  x_i^*(M,p_m)
  =
  \left[
  \frac{\beta a_i}{\kappa}\Omega_i(M,p_m)
  \right]^{1/\nu},
  \label{eq:p16m-advancement}
\end{equation}
including \(x_i^*=0\) when \(\Omega_i=0\). The binary reform affects
\(x_i^*\) only through \(\Omega_i\). At a fixed \(p_m\), its effect is
strict only if \(E\) improves on the pre-MAH maximum on a
positive-probability project set.

\subsection{Qualified manufacturing-service equilibrium}
\label{subsec:p16m-cmo}

Qualified supplier \(j\), with efficiency \(z_j\), chooses capacity
\(s_j\geq0\):
\begin{equation}
  \max_{s_j\geq0}\{p_ms_j-\Psi(s_j;z_j)\}.
  \label{eq:p16m-supplier}
\end{equation}
With \(\Psi_{ss}>0\), positive-price capacity is uniquely determined by
\(p_m=\Psi_s(s_j^*;z_j)\), and aggregate supply is
\[
  S_m(p_m)=\int s_j^*(p_m,z)\,dH_C(z).
\]

Expected entrusted capacity per planning-stage project is
\begin{equation}
  \chi_i^E(p_m;M)
  =
  \int b(m)\mathbf 1\{r_i^*(q,m;M,p_m)=E\}\,dF(q,m).
  \label{eq:p16m-project-demand}
\end{equation}
Study-related and total demand are
\begin{equation}
  \begin{aligned}
  D_m^{\mathrm{MAH}}(p_m;M)
  &=
  \int a_ix_i^*(M,p_m)\chi_i^E(p_m;M)\,dH(a,k),\\
  D_m(p_m;M)
  &=D_m^B(p_m)+D_m^{\mathrm{MAH}}(p_m;M).
  \end{aligned}
  \label{eq:p16m-demand}
\end{equation}
The demand expression includes both feedbacks: price changes expected value
and hence advancement, and it changes deterministic route selection.
Continuous heterogeneity makes aggregate demand continuous even though
individual indicators can jump.

The qualified-capacity market clears at the unique price
\begin{equation}
  D_m(p_m^*;M)=S_m(p_m^*).
  \label{eq:p16m-clearing}
\end{equation}
Supply is strictly increasing and demand weakly decreasing. Positive
background demand gives excess demand at zero, while entrusted and
background demand vanish and supply diverges at high prices. These conditions
give existence and uniqueness. If post-MAH study demand is positive at the
old price, \(p_m^*(1)>p_m^*(0)\): CMO scarcity attenuates the reform rather
than amplifying it through a direct price reduction.

The baseline partial equilibrium contains exactly
\begin{equation}
  \left\{
  p_m^*,\;
  x_i^*,\;
  r_i^*(q,m),\;
  s_j^*
  \right\}_{i,j}.
  \label{eq:p16m-equilibrium}
\end{equation}
No entry, labor, capital, product-market aggregate, or welfare closure is
added.

\subsection{Main predictions}
\label{subsec:p16m-predictions}

\paragraph{Organizational sorting.}
Conditional on \(I\) and \(E\) dominating the outside options, low internal
manufacturing capability selects entrusted production and high capability
selects internal production. Lower entrusted-route friction expands the
entrusted region; a higher fixed CMO price contracts it.

\paragraph{MAH-relevant projects and advancement.}
The reform has a zero effect whenever the newly feasible \(E\) value does not
exceed the best old route. Otherwise it raises expected project value and,
at fixed service price, raises common project advancement. The response is
larger where the entrusted value gain is larger; MAH does not change
\(a_i\), project quality, or the advancement-cost function.

\paragraph{CMO equilibrium and scarcity.}
The qualified-capacity price is unique. Reform-induced entrusted demand can
raise this price, weakly reducing the direct fixed-price gains in project
value and advancement. Fixed-price comparisons and equilibrium-price
comparisons are therefore distinct.

\paragraph{Planning, organization, and observed outcomes.}
At equilibrium,
\[
  \Lambda_i^{\mathrm{plan}}=a_ix_i^*,
\]
while expected realized retained outcomes are
\begin{equation}
  Y_i^{\mathrm{ret}}
  =
  a_ix_i^*
  \int s(q)
  \left[
  \mathbf 1\{r_i^*=I\}
  +
  \mathbf 1\{r_i^*=E\}
  \right]dF(q,m).
  \label{eq:p16m-retained-outcome}
\end{equation}
Thus observed outcomes can change through advancement or route composition.
Holder--producer separation is observed only after route assignment and
cannot cause the earlier \(x_i\) choice.

\paragraph{Novelty composition.}
For \(g\in\{O,\mathrm{Inc}\}\), outcomes use the same \(x_i^*\) and the
mixture \(\rho_gF_g\). The baseline permits either class to have a zero or
larger response and imposes no original-versus-incremental ranking. Patent
applications are not a baseline endogenous outcome.

\end{document}
